\caption{\textbf{How correlated are neurons' responses, and is there structure beyond these correlations?}
\textbf{a} Cross-validated correlation matrix of the feature vectors of all 54,569 neurons with both trial halves. Entry $(i,j)$ is the correlation between neuron $i$'s odd-trial and neuron $j$'s even-trial feature vector, symmetrised, so only correlations that replicate across trial halves remain. The diagonal is each neuron's odd/even reliability (median 0.27). Neurons are in the Rastermap order fitted on odd trials (Fig.~\ref{fig:latency_tiling}a,b; ticks, cluster numbers). Bars: dense (orange) and sparse tiling (blue). The matrix is shown as exact means over blocks of $27\times27$ neurons. Two large, mutually anticorrelated groups dominate: clusters $\sim$13--35 (elevated before movement onset, suppressed at reward) and $\sim$78--99 (feedback). Sparse-tiling neurons correlate weakly with all others. Mean $|r|$ over all $3.0\times10^9$ pairs: 0.11.
\textbf{b} Dense-tiling neurons (box in \textbf{a}; 4,251) at full resolution; lines, cluster boundaries. Correlation is highest within a cluster and falls off with distance along the latency tiling.
\textbf{c--e} Structure beyond pairwise correlations. Data: all 54,719 neurons, all trials, top 10 (\textbf{c}, \textbf{d}) or 20 principal components of the feature vectors (\textbf{e}, left). All nulls have the data's covariance:
Gaussian (blue);
PC-marginal, principal-component scores permuted independently (orange);
Gaussian copula, the exact distribution of each score plus their rank correlations (purple);
ICA-marginal, independent components (FastICA) permuted independently, i.e.\ random mixing of independent sparse components (green).
Positive control: a 25-cluster Gaussian mixture fitted to the data (pink). 10 null draws each.
\textbf{c} ePAIRS-style test \cite{raposo2014category,hirokawa2019frontal}: each neuron's mean angle to its 5 nearest neighbours, with neighbours taken from other probe insertions only. AUC is the probability that a real neuron's angle is smaller than a null neuron's (0.5, no excess clustering), here against the Gaussian / PC-marginal nulls.
\textbf{d} k-means silhouette against the number of clusters. Real neurons are slightly more clusterable than the nulls but have no preferred number of clusters, unlike the mixture control (bump near 20).
\textbf{e} AUC against each null for response profiles (left) and for task-variable selectivity (right): RRR coefficients for 8 task variables summed over time, signed or absolute (38,853 neurons from 169 sessions; neighbours from other sessions; Methods; the features of Fig.~\ref{fig:rrr_peth_selectivity_comparison}). Response profiles are more clustered than every null (ICA-marginal: 0.55 and 0.59), about as much as the mixture control. Task-variable selectivity is close to random mixing (ICA-marginal: 0.51 signed, 0.52 absolute), although the test detects a 25-cluster mixture in the same space (0.56). Absolute values are non-negative and skewed, which only the copula and ICA-marginal nulls reproduce.}
\label{fig:response_correlations}
